\exercisesection{重数}

\begin{enumerate}
\def\labelenumi{\arabic{enumi})}
\tightlist
\item
  设 \(\Gamma \subset \mathbf{N}\) 为 \(\mathbf{N}\) 的子幺半群，使得 \(\mathbf{N} - \Gamma\) 是有限集。设 \((a_1, \ldots, a_r)\)（其中 \(0 < a_1 < \cdots < a_r\)）为 \(\Gamma\) 的极小生成元系，且 \(k\) 为一个域。置 \(\mathbf{A} = k[\Gamma]\)。
\end{enumerate}

\begin{enumerate}
\def\labelenumi{\alph{enumi})}
\item
  证明 A 同构于 \(k[\mathbf{T}]\) 中由 \(\mathbf{T}^{a_1}\)、\(\ldots\)、\(\mathbf{T}^{a_r}\) 生成的 k-子代数。
\item
  设 m 为由 \(T^{a_{1}}\)、\(\ldots\)、\(T^{a_{r}}\) 生成的 A 的极大理想；设 \(A_{m}\) 为 A 在 m 处的局部环，并置 \(q = mA_{m}\)。证明有 \(\dim_{k}q/q^{2} = r\) 及 \(e_{q}(A_{m}) = a_{1}\)。
\end{enumerate}

\begin{enumerate}
\def\labelenumi{\arabic{enumi})}
\setcounter{enumi}{1}
\tightlist
\item
  设 \(k\) 为一个域，\(m\) 为满足 \(\geqslant 1\) 的整数，I 为多项式环 \(k[\mathbf{T}_{1},\ldots,\mathbf{T}_{m}]\) 的一个理想，它由 \(\mathbf{M}_{1},\ldots,\mathbf{M}_{s}\) 所表示的单项式生成，这些点属于 \(\mathbf{R}_{+}^{m}\)，即 \(\mathbf{R}^{m}\) 的“正象限”。记 \(\mathcal{N}\) 为 \(\mathbf{R}_{+}^{m}\) 中的集合 \(\mathbf{M}_{i}+\mathbf{R}_{+}^{m}\)（\(1\leqslant i\leqslant s\)）之并的凸包。
\end{enumerate}

\begin{enumerate}
\def\labelenumi{\alph{enumi})}
\tightlist
\item
  置 \(m = (T_1, \ldots, T_m)\)、\(A = k[T_1, \ldots, T_m]_m\) 及 \(q = I.A\)。证明 \(A/q\) 具有有限长度，当且仅当 \(\mathbf{R}_+^m - \mathcal{N}\) 的体积有限。
\item
  证明等式 \(e_q(A) = m! \operatorname{Vol}(\mathbf{R}_+^m - \mathcal{N})\)。
\end{enumerate}

\begin{enumerate}
\def\labelenumi{\arabic{enumi})}
\setcounter{enumi}{2}
\tightlist
\item
  设 \(k\) 为一个域，\(f_{1},...,f_{r}\) 为分次环 \(k[X_{1},...,X_{n}]\) 中构成完全相交序列的齐次元素。称商环
\end{enumerate}

\[
\mathbf {H} = k [ \mathrm{X} _ {1},..., \mathrm{X} _ {n} ] / (f _ {1},..., f _ {r})
\]

  是分次完全相交环。设 I 为环 \(\mathbf{H} = k[\mathbf{X}_1, \ldots, \mathbf{X}_n] / (f_1, \ldots, f_r)\) 的非零分次理想，并赋予其商分次。证明有 \(\dim(\mathbf{I}) = \dim(\mathbf{H})\)。推出或者 \(\dim(\mathbf{H}/\mathbf{I}) < \dim(\mathbf{H})\)，或者 \(\dim(\mathbf{H}/\mathbf{I}) = \dim(\mathbf{H})\) 且 \(c_{\mathbf{M}/\mathbf{I}} < c_{\mathbf{M}}\)（\(\S 4, \text{n}^{\circ} 2\)）。（考虑由 \(s \in \mathbf{H}\) 中满足 \(s.\mathbf{I} = \{0\}\) 的元素构成的理想 b，并证明若 \(\dim(\mathbf{I}) < \dim(\mathbf{H})\)，则 b 含有一个非零因子元素；证明方法是考察 b 相对于 H 的极小素理想及 \(\mathbf{H}_{+} = \mathbf{H}_{\geqslant 1}\) 的位置。））

\begin{enumerate}
\def\labelenumi{\arabic{enumi})}
\setcounter{enumi}{3}
\tightlist
\item
  设 A 为诺特局部环，极大理想为 m。
\end{enumerate}

\begin{enumerate}
\def\labelenumi{\alph{enumi})}
\item
  假设 \(\mathrm{gr}_{\mathfrak{m}}(\mathrm{A})\) 是分次完全相交环（参见习题 3）。证明 A 的完备化 \(\widehat{A}\) 是完全相交的，即是正则局部环对一个由完全相交序列生成的理想的商。
\item
  设 \(x \in \mathfrak{m}\)，且 \(\xi\) 为 \(x\) 在 \(\mathfrak{m}^{\nu}/\mathfrak{m}^{\nu+1}\) 中的像，其中 \(\nu\) 满足 \(x \in \mathfrak{m}^{\nu}\) 且 \(x \notin \mathfrak{m}^{\nu+1}\)。证明 \(e_{\mathfrak{m}}(\mathbf{A}/x\mathbf{A}) = \nu e_{\mathfrak{m}}(\mathbf{A})\) 当且仅当 \(\xi\) 不是 \(\mathrm{gr}_{\mathfrak{m}}(\mathbf{A})\) 中的零因子，并且在此情形下有 \(\mathrm{gr}_{\mathfrak{m}}(\mathbf{A}/x\mathbf{A}) = \mathrm{gr}_{\mathfrak{m}}(\mathbf{A})/\xi\mathrm{gr}_{\mathfrak{m}}(\mathbf{A})\)，且 \(x\) 不是 A 中的零因子。（令 I 为 \(\mathrm{gr}_{\mathfrak{m}}(\mathbf{A})\) 的齐次理想，由 \(\eta \in \mathrm{gr}_{\mathfrak{m}}(\mathbf{A})\) 中满足 \(\xi\eta = 0\) 的元素构成。考虑 A 的理想 \(\mathrm{N}_n = \mathfrak{m}^{\nu+n}:x\)，以及从 \(\mathrm{I}_n\) 到 \(\mathrm{N}_{n+1}/\mathfrak{m}^{n+1}\) 的映射，它将 \(\eta \in \mathrm{I}_n\) 送到模 \(\mathfrak{m}^{n+1}\) 下的、元素 \(y \in A\) 提升 \(\eta\) 后的类。我们将证明它是单射，并由此推出不等式
\end{enumerate}

\[
\operatorname{long} _ {\mathbf {A} / \mathfrak {m}} (\mathbf {I} _ {n}) \leqslant \operatorname{long} _ {\mathbf {A} / \mathfrak {m}} (\mathbf {N} _ {n + 1} / \mathfrak {m} ^ {n + 1});\tag{i}
\]

此外，我们将证明等式

\[
(1 - \mathrm{T} ^ {\mathrm{v}}) \mathrm{H} _ {\mathrm{A}} ^ {(1)} = \mathrm{T} ^ {- \mathrm{v}} \mathrm{H} _ {\mathrm{A} / x \mathrm{A}} ^ {(1)} - \sum_ {n \geqslant \mathrm{v}} \operatorname{long} _ {\mathrm{A} / \mathfrak {m}} (\mathrm{N} _ {n} / \mathfrak {m} ^ {n}). \mathrm{T} ^ {n}.\tag{ii}
\]

最后注意到，\(\xi\) 是零因子当且仅当 I \(\neq\) \{0\}；利用上一习题，由 (i) 推出 \(\sum_{n\geqslant 0}\mathrm{long}_{\mathrm{A / m}}(\mathbf{N}_{n + 1} / \mathfrak{m}^{n + 1})\) \(\mathbf{T}^n = \frac{\mathbf{S}(\mathbf{T})}{(1 - \mathbf{T})^\nu}\)，其中 \(\mathbf{S}(1)\neq 0\)，并由 (ii) 得到不等式 \(e_{\mathfrak{m}}(\mathbf{A} / x\mathbf{A}) > e_{\mathfrak{m}}(\mathbf{A})\)。c) 设 \((x_{1},\dots,x_{s})\) 为极大理想 m 中元素组成的序列。假设 \(x_{i}\in \mathfrak{m}^{\nu_i}\) 且 \(x_{i}\notin \mathfrak{m}^{\nu_i + 1}\)，并令 \(\xi_{i}\) 为 \(x_{i}\) 在 \(\mathfrak{m}^{\nu_i} / \mathfrak{m}^{\nu_i + 1}\) 中的类。证明有 \(e_{\mathfrak{m}}(\mathbf{A} / x)\geqslant e_{\mathfrak{m}}(\mathbf{A})\nu_{1}\dots \nu_{s}\)，且等号成立当且仅当序列 \((\xi_1,\dots,\xi_s)\) 在环 \(\mathrm{gr}_{\mathfrak{m}}(\mathbf{A})\) 中是完全相交的。在此情形下，\(\mathrm{gr}_{\mathfrak{m}}(\mathbf{A} / x)\) 与 \((\mathrm{gr}_{\mathfrak{m}}(\mathbf{A}))/\xi\) 之间有分次环同构，其中 \(\mathbf{x} = \sum_{i}x_{i}\mathbf{A}\)，\(\xi = \sum_{i}\xi_{i}\mathrm{gr}_{\mathfrak{m}}(\mathbf{A})\)。

\begin{enumerate}
\def\labelenumi{\arabic{enumi})}
\setcounter{enumi}{4}
\tightlist
\item
  设 \(k\) 为一个域。置 \(\mathbf{R} = k[[\mathrm{X},\mathrm{Y},\mathrm{Z}]]\)，并考虑极大理想为 m 的局部环 \(\mathbf{A} = \mathbf{R}/(\mathbf{XY},\mathbf{X}\mathbf{Z})\)。证明有
\end{enumerate}

\[
\mathrm{H} _ {\mathrm{A}} = (1 - \mathrm{T}) (1 - \mathrm{T} - \mathrm{T} ^ {2}) \mathrm{H} _ {\mathrm{R}} = 1 + \sum_ {n \geqslant 1} (n + 2) \mathrm{T} ^ {n}.
\]

因此有 \(\dim(A) = 2\)、\(\dim_k(m/m^2) = 3\) 及 \(e_m(A) = 1\)。（我们有 \(H_A = P_G\)，其中 \(G = gr_m(A) = k[X, Y, Z]/(XY, XZ)\)。我们将证明存在分次 H-模的正合序列

\[
0 \rightarrow \mathrm{H} (- 3) \rightarrow \mathrm{H} (- 2) \oplus \mathrm{H} (- 2) \rightarrow \mathrm{H} \rightarrow \mathrm{G} \rightarrow 0
\]

\(\text{其中 } \mathrm{H} = k[\mathrm{X},\mathrm{Y},\mathrm{Z}].\)

\begin{enumerate}
\def\labelenumi{\arabic{enumi})}
\setcounter{enumi}{5}
\item
  设 \(k\) 为一个域，A 为维数 1 的局部、完备、约化 \(k\)-代数，\(m\) 为其极大理想。证明 A 的一个理想的最小生成元数至多为 \(e_m(A)\)。存在 \(X \in A\)，使得从 \(k[[X]]\) 到 A 的单射使 A 成为 \(k[[X]]\)-自由模，秩为 \(e_m(A)\)。
\item
  设 A 为极大理想为 m 的局部环，M 为有限生成 A-模，\((x_{1}, \ldots, x_{d})\) 为 m 中元素组成的、对 A-模 M 相交的序列，x 为理想 \(x_{1}A + \cdots + x_{d}A\)。
\end{enumerate}

\begin{enumerate}
\def\labelenumi{\alph{enumi})}
\tightlist
\item
  假设 A-模 M 是忠实的。则 \(\mathrm{gr}_{x}(\mathbf{A})\) 除以其幂零根基 n 所得的商同构于 \((\mathbf{A}/\mathfrak{m})[\mathbf{X}_{1}, \ldots, \mathbf{X}_{d}]\)，且 \((\mathrm{gr}_{x}(\mathbf{M}))_{n}\) 是一个 \((\mathrm{gr}_{x}(\mathbf{A}))_{n}\)-模，有限长度为 \(e_{x}(\mathbf{M})\)。
\item
  一般情形下，\(\mathrm{gr}_{x}(\mathbf{A})\) 有唯一的极小素理想 n，且 \((\mathrm{gr}_{x}(\mathbf{M}))_{n}\) 是一个 \(\mathrm{gr}_{x}(\mathbf{A})\)-模，有限长度为 \(e_{x}(\mathbf{M})\)。
\end{enumerate}

\begin{enumerate}
\def\labelenumi{\arabic{enumi})}
\setcounter{enumi}{7}
\tightlist
\item
  设 A 为一个环，a 为 A 的一个理想。置 \(a^{i} = A\) 对 \(i \leqslant 0\) 成立，并考虑 A-代数 \(\mathcal{A}(\mathfrak{a}) = \sum_{i \in \mathbf{Z}} \mathfrak{a}^{-i} v^{i} \subset \mathrm{A}[v, v^{-1}]\) 及 A-代数 \(\mathrm{P}(\mathfrak{a}) = \sum_{i \geqslant 0} \mathfrak{a}^{i} v^{i} \subset \mathrm{A}[v]\)。这两个分次 A-代数具有相同的全分式环。证明对于元素 \(a \in A\)，下列条件等价：
\end{enumerate}

\begin{enumerate}
\def\labelenumi{(\roman{enumi})}
\item
  元素 \(av\) 属于 \(\mathcal{A}(\mathfrak{a})\)-代数 \(\mathbf{A}[v, v^{-1}]\)，且它在 \(\mathcal{A}(\mathfrak{a})\) 上整。
\item
  P(a)-代数 A{[}v{]} 的元素 av 在 P(a) 上整。
\item
  存在整依赖关系：
\end{enumerate}

\[
a ^ {k} + b _ {1} a ^ {k - 1} + \dots + b _ {k} = 0 \quad \text { with } \quad b _ {j} \in \mathfrak {a} ^ {j}.
\]

称 \(a\) 关于理想 \(\mathfrak{a}\) 整。

\begin{enumerate}
\def\labelenumi{\arabic{enumi})}
\setcounter{enumi}{8}
\tightlist
\item
  \begin{enumerate}
  \def\labelenumii{\alph{enumii})}
  \tightlist
  \item
    设 A 和 a 如上，b 为 A 的一个包含 a 的理想。证明下列条件等价：
  \end{enumerate}
\end{enumerate}

\begin{enumerate}
\def\labelenumi{(\roman{enumi})}
\item
  b 的每个元素都关于 a 整。
\item
  \(\mathcal{A}(\mathfrak{a})\)-分次代数 \(\mathcal{A}(\mathfrak{b})\) 是整的。
\item
  分次 \(P(\mathfrak{a})\)-代数 \(P(\mathfrak{b})\) 是整的。
\item
  存在整数 \(k \geqslant 0\)，使得 \(\mathfrak{a}\mathfrak{b}^{k} = \mathfrak{b}^{k+1}\)，从而有 \(\mathfrak{a}^{n}\mathfrak{b}^{k} = \mathfrak{b}^{k+n}\) 对所有正整数 \(n\) 成立。称 \(\mathfrak{b}\) 关于 \(\mathfrak{a}\) 整。
\end{enumerate}

\begin{enumerate}
\def\labelenumi{\alph{enumi})}
\setcounter{enumi}{1}
\item
  证明在关于 a 整的理想中存在最大的理想。记其为 \(\overline{a}\)，称为 a 在 A 中的整闭包。验证若 \(a \subset b\)，则 \(\overline{a} \subset \overline{b}\)，且 \(\overline{\overline{a}} = \overline{a}\)。证明若整闭包 \(\overline{P(a)}\) 是有限生成的 \(P(a)\)-模，则当 k 充分大时，\(\mathfrak{a}\,\overline{\mathfrak{a}^{k}} = \overline{\mathfrak{a}^{k+1}}\)。
\item
  证明若 \(a \in A\) 满足整依赖关系
\end{enumerate}

\[
a ^ {k} + b _ {1} a ^ {k - 1} + \dots + b _ {k} = 0 \quad \text { with } \quad k \geqslant 1 \quad \text { and } \quad b _ {i} \in \mathfrak {a} ^ {i},
\]

 则 \(\mathfrak{a}(\mathfrak{a} + \mathbf{A}a)^{k-1} = (\mathfrak{a} + \mathbf{A}a)^k\)。由此推出，\(\mathfrak{a}\) 的整闭包在 A 中恰好是 \(\mathbf{A}\) 中关于 \(\mathfrak{a}\) 整的元素集合。

10) 设 A 为一个环。

\begin{enumerate}
\def\labelenumi{\alph{enumi})}
\item
  设 B 为整的 A-代数。证明对于 A 的每个元素 a 及每个理想 \(\mathfrak{a}\)，元素 a 关于 \(\mathfrak{a}\) 整，当且仅当 a 关于 B 中理想 \(\mathfrak{a}B\) 整，该理想由 \(\mathfrak{a}\) 生成。
\item
  假设 A 是诺特约化环，记 \(\mathfrak{p}_{1},\ldots,\mathfrak{p}_{r}\) 为 A 的极小素理想，并考虑自然单射：
\end{enumerate}

\[
\mathrm{A} \rightarrow \prod_ {1 \leqslant i \leqslant r} \mathrm{A} / \mathfrak {p} _ {i}.
\]

证明它使 \(\prod_{1\leqslant i\leqslant r}\mathrm{A} / \mathfrak{p}_i\) 成为整的 A-代数。

\begin{enumerate}
\def\labelenumi{\alph{enumi})}
\setcounter{enumi}{2}
\item
  由此推出，给定诺特环 A，A 的元素 a 关于理想 \(\mathfrak{a}\) 整的充要条件是：对于 A 的每个极小素理想 p，a 在 A/p 中的像关于理想 \((\mathfrak{a} + p)/p\) 整。
\item
  证明若 \(\overline{a} = \overline{b}\)，则理想 \(a\) 和 \(b\) 具有相同的极小相伴素理想。
\end{enumerate}

\begin{enumerate}
\def\labelenumi{\arabic{enumi})}
\setcounter{enumi}{10}
\tightlist
\item
  设 A 为诺特环，a、b 为 A 的两个理想，且 Supp(a) = Supp(b) = Spec(A)。证明下列条件等价：
\end{enumerate}

\begin{enumerate}
\def\labelenumi{(\roman{enumi})}
\tightlist
\item
  存在有限生成的 A-模 M，使得 \(\operatorname{Supp}(\mathbf{M}) = \operatorname{Spec}(\mathbf{A})\)，且 \(bM \subset aM\)。\\
\item
  有 \(a, b \subset \overline{a}\)。
\end{enumerate}

\begin{enumerate}
\def\labelenumi{\arabic{enumi})}
\setcounter{enumi}{11}
\tightlist
\item
  设 A 为诺特环，a、b、c 为 A 的三个理想，且 c 包含于 A 的根基中。
\end{enumerate}

\begin{enumerate}
\def\labelenumi{\alph{enumi})}
\item
  证明若有包含关系 \(\overline{b + c.a} \subset \overline{a}\)，则有 \(\overline{b} \subset \overline{a}\)。
\item
  因此，若 A 是局部环，且有 \(b \subset a\) 和 \(\overline{b} = \overline{a}\)，则至少存在一个理想 \(b_{1} \subset b\)，使得
\end{enumerate}

\[
\overline {{{\mathbf {b} _ {1}}}} = \overline {{{\mathbf {a}}}},
\]

\(\beta)\) \(\mathbf{b}_{1}\) 对此性质是极小的，即若 \(\mathbf{b}_{1}^{\prime}\subset \mathbf{b}_{1}\) 且 \(\overline{\mathbf{b}_1'} = \overline{\mathbf{b}_1}\)，则 \(\mathbf{b}_1' = \mathbf{b}_1\)。

\begin{enumerate}
\def\labelenumi{\arabic{enumi})}
\setcounter{enumi}{12}
\item
  设 A 为局部环，a、b 为 A 中异于 A 的两个理想。证明若 a ⊂ b，则 \(\overline{a} = \overline{b}\) 当且仅当分次代数 gr \(_b(A)\) 在分次代数 gr \(_a(A)\) 上整。推出若 a = (x \(_1\) , \ldots, x \(_d\) ) ⊂ b，其中 d = dim(A)，则 \(\overline{a} = \overline{b}\) 当且仅当初始形式 \(\xi_{1}\) , \ldots, \(\xi_{d}\)（它们是在分次代数 gr \(_b(A)\) 中的 x \(_i\) 的初始形式）次数为 1，且构成 gr \(_b(A)\) 的极大相交序列。
\item
  设 A 为局部环，m 为其极大理想，q 为 A 中异于 A 的理想，且 A/q 具有有限长度。假设 A 含有子域 k，使得 A = k + m。
\end{enumerate}

\begin{enumerate}
\def\labelenumi{\alph{enumi})}
\tightlist
\item
  证明 q 中存在元素 \(x_{1}, \ldots, x_{d}\)，其中 \(d = \dim(A)\)，且 \(x_{i} \in q^{\delta_{i}}\)，使得对于每个充分大的整数 v，有 \(q^{v} = x_{1} \cdot q^{v - \delta_{1}} + \cdots + x_{d} \cdot q^{v - \delta_{d}}\)。
\item
  若假设域 A/m 是无限的，则可以找到一个由 q 中 d 个元素生成的理想 a，其中 \(d = \dim(A)\)，使得 \(\overline{a} = \overline{q}\)（可以使用 § 7, no. 5。）
\end{enumerate}

\begin{enumerate}
\def\labelenumi{\arabic{enumi})}
\setcounter{enumi}{14}
\tightlist
\item
  设 A 为极大理想 m 的局部、诺特、完备环。置 \(d = \dim(A)\)。假设 A 含有子域 k，使得 \(A = k + m\)。
\end{enumerate}

证明可以在 m 中找到 \(x_{1},\ldots,x_{d}\)，使得：

  态射 \(\varphi:k[[X_{1},...,X_{d}]]\to A\) 由 \(\varphi(X_{i})=x_{i}\) 定义，使 A 成为整的 \(k[[X_{1},...,X_{d}]]\)-代数；

\(\beta)\) 记 \(\mathfrak{m}_0 = (X_1, \ldots, X_d)\) 为 \(k[[X_1, \ldots, X_d]]\) 的极大理想，则有 \(\overline{\mathfrak{m}_0A} = \mathfrak{m}\)，因而 \(e_{\mathfrak{m}_0\mathbf{A}}(\mathbf{A}) = e_{\mathfrak{m}}(\mathbf{A})\)。

  证明若 \(x_{1},\ldots,x_{d}\) 满足条件 \(\alpha\)，则条件 \(\beta\) 成立当且仅当初始形式 \(\xi_{1},\ldots,\xi_{d}\)（它们是 \(x_{i}\) 在 \(\mathrm{gr}_{\mathfrak{m}}(\mathbf{A})\) 中的初始形式）构成极大相交序列。

\begin{enumerate}
\def\labelenumi{\arabic{enumi})}
\setcounter{enumi}{15}
\item
  设 A 为诺特局部环，\(x = (x_1, \ldots, x_d)\) 为由 A 的极大相交序列生成的理想。证明有不等式 \(e_x(A) \leqslant \text{long}(A/x)\)，且等号成立只能在同态 \((A/x)[T_1, \ldots, T_d] \xrightarrow{\varphi} \text{gr}_x(A)\) 为同构时成立；该同态由 \(\varphi(T_i) = \xi_i\) 定义，其中 \(\xi_i\) 是 \(x_i\) 模 \(x^2\) 的类（这意味着 \((x_1, \ldots, x_d)\) 是完全相交的）。
\item
  设 A 为诺特局部环，具有长度 \(d = \dim(A)\) 的完全相交序列；令 \(\mathfrak{a}\) 为由此序列生成的理想，令 \(f \in A\) 为关于 \(\mathfrak{a}\) 整但不属于该理想的元素（例如 \(A = k[[X_1, ..., X_d]]\)，其中 k 为域，\(\mathfrak{a} = (X_1^k, ..., X_d^k)\)，且 \(f = X_1^{k-1}X_2\)）。令 \(\mathfrak{a}' = \mathfrak{a} + fA\)。证明有不等式 \(e_{\mathfrak{a}'}(A) > \text{long}(A/\mathfrak{a}')\)。
\item
  设 A 为一个环，\(\mathfrak{a}\) 为 A 的一个理想。对于每个元素 \(a \in \mathbf{A}\)，以 \(\nu_{\mathfrak{a}}(a)\) 表示整数 \(\nu \in \mathbf{N}\) 中满足 \(a \in \mathfrak{a}^{\nu}\) 者的集合的上界。
\end{enumerate}

\begin{enumerate}
\def\labelenumi{\alph{enumi})}
\tightlist
\item
  证明若 \(a \notin \bigcap_{k=0}^{\infty} \mathfrak{a}^k\)，则序列 \(\left(\frac{\nu_{\mathfrak{a}}(a^k)}{k}\right)\) 收敛。以 \(\overline{\nu}_{\mathfrak{a}}(a)\) 表示其极限。若
\end{enumerate}

\(a \in \bigcap_{k=0}^{\infty} \mathfrak{a}^{k}\)，则置 \(\overline{\nu}_{\mathfrak{a}}(a) = +\infty\)。

\begin{enumerate}
\def\labelenumi{\alph{enumi})}
\setcounter{enumi}{1}
\tightlist
\item
  证明对于所有 \(a \in \mathbf{A}\)，有 \(\overline{\nu}_{\mathfrak{a}}(a) = \overline{\nu}_{\mathfrak{a}}(a)\)，并由此推出下列条件等价：
\end{enumerate}

\begin{enumerate}
\def\labelenumi{(\roman{enumi})}
\item
  \(\overline{\nu}_{\mathfrak{a}}(a)\geqslant 1,\)
\item
  \(a \in \overline{\mathfrak{a}}\).
\end{enumerate}

（研究 a) 中所考察序列的单调性。）

\begin{enumerate}
\def\labelenumi{\arabic{enumi})}
\setcounter{enumi}{18}
\tightlist
\item
  设 A 为诺特环，q 为 A 的一个理想，且 A/q 具有有限长度，q 包含于 A 的根基中。证明有
\end{enumerate}

\[
e _ {\mathfrak {q}} (\mathrm{A}) = e _ {\overline {{\mathfrak {q}}}} (\mathrm{A}).
\]

\begin{enumerate}
\def\labelenumi{\arabic{enumi})}
\setcounter{enumi}{19}
\tightlist
\item
  设 A 为诺特环，M 为有限生成的 A-模，\(q_{1}, \ldots, q_{k}\) 为 A 的理想，且对于每个 i，\(M/q_{i}M\) 具有有限长度。假设 \(M \neq 0\)，并置 \(d = \dim_{\mathbf{A}}(\mathbf{M})\)。a) 证明 \(M/q_{1}^{n_{1}} \ldots q_{k}^{n_{k}}M\) 对每个由正整数组成的 k-元组 \((n_{1}, \ldots, n_{k})\) 都具有有限长度。
\end{enumerate}

\begin{enumerate}
\def\labelenumi{\alph{enumi})}
\setcounter{enumi}{1}
\tightlist
\item
  考虑类型为 \(\mathbf{N}^k\) 的分次环 H，使得
\end{enumerate}

\[
\mathrm{H} _ {n _ {1}, \dots , n _ {k}} = \mathfrak {q} _ {1} ^ {n _ {1}} \mathfrak {q} _ {2} ^ {n _ {2}} \dots \mathfrak {q} _ {k} ^ {n _ {k}} / \mathfrak {q} _ {1} ^ {n _ {1} + 1} \mathfrak {q} _ {2} ^ {n _ {2}} \dots \mathfrak {q} _ {k} ^ {n _ {k}},
\]

以及分次 H-模 N，使得

\[
 \mathrm{N} _ {n _ {1}, \dots , n _ {k}} = \mathrm{q} _ {1} ^ {n _ {1}} \mathrm{q} _ {2} ^ {n _ {2}} \dots \mathrm{q} _ {k} ^ {n _ {k}} \mathrm{M} / \mathrm{q} _ {1} ^ {n _ {1} + 1} \mathrm{q} _ {2} ^ {n _ {2}} \dots \mathrm{q} _ {k} ^ {n _ {k}} \mathrm{M};
\]

证明 N 是有限生成的分次 H-模，由 \(N_{0,\ldots,0}\) 及有限个次数为 \(\mathbf{Z}^{k}\) 的标准基向量的元素生成。

由此推出，存在多项式 \(\mathbf{Q}(\mathrm{T}) \in \mathbf{Z}[\mathrm{T}_1,..,\mathrm{T}_k]\) 及整数 \(s_1 \geqslant 1, ..., s_k \geqslant 1\)，使得有

\[
\sum_ {n \in \mathbb {N} ^ {k}} \operatorname{long} _ {\mathrm{A}} (\mathrm{H} _ {n}). \mathrm{T} ^ {n} = \frac {\mathrm{Q} (\mathrm{T})}{(1 - \mathrm{T} _ {1}) ^ {s _ {1}} \dots (1 - \mathrm{T} _ {k}) ^ {s _ {k}}}.
\]

（我们对 d 使用归纳法。）

\begin{enumerate}
\def\labelenumi{\alph{enumi})}
\setcounter{enumi}{2}
\tightlist
\item
  由此，通过将 § 6, no. 1 的结果延拓到 \(\mathbf{Z}((\mathrm{T}_{1},\ldots,\mathrm{T}_{k}))\)，推出存在整数 \(c(\alpha_{1},\ldots,\alpha_{k})\)，对应于序列 \((\alpha_{1},\ldots,\alpha_{k})\)，其中 \(|\alpha| = \sum_{i=1}^{k}\alpha_{i} = d\)，使得有
\end{enumerate}

\[
e _ {\mathfrak {q} _ {1} ^ {\nu_ {1}} \dots \mathfrak {q} _ {k} ^ {\nu_ {k}}} (\mathbf {M}) = \sum_ {| \alpha | = d} \frac {d !}{\alpha_ {1} ! \dots \alpha_ {k} !} c (\alpha_ {1},..., \alpha_ {k}) v _ {1} ^ {\alpha_ {1}}... v _ {k} ^ {\alpha_ {k}}
\]

其中 \((\mathbf{v}_1,\dots,\mathbf{v}_k)\) 属于 \(\mathbf{N}^k\)

\begin{enumerate}
\def\labelenumi{\alph{enumi})}
\setcounter{enumi}{3}
\tightlist
\item
  现在假设 A 是剩余域为无限域 κ 的局部环。利用分次 H-模 N 的表面元素的存在性，证明有
\end{enumerate}

\[
c (\alpha_ {1}, \dots , \alpha_ {k}) = \inf _ {\mathbf {q} _ {\alpha}} (e _ {\mathbf {q} _ {\alpha}} (\mathbf {M}))
\]

  其中 \(q_{\alpha}\) 遍历 A 的理想集，这些理想分别由 \(\alpha_{1}\) 个 \(q_{1}\) 中的元素、\(\alpha_{2}\) 个 \(q_{2}\) 中的元素、\ldots、\(\alpha_{k}\) 个 \(q_{k}\) 中的元素生成。有时我们将 \(e(\mathbf{q}_{1}^{[\alpha_{1}]} + \cdots + \mathbf{q}_{k}^{[\alpha_{k}]}; \mathbf{M})\) 记为整数 \(c(\alpha_{1}, \ldots, \alpha_{k})\)；当 M = A 时，简写为 \(e(\mathbf{q}_{1}^{[\alpha_{1}]} + \cdots + \mathbf{q}_{k}^{[\alpha_{k}]})\)。e) 证明等式

\[
e _ {\mathfrak {q} _ {1}, \mathfrak {q} _ {2}} (\mathbf {M}) = \sum_ {i = 1} ^ {d} \binom {d} {i} e (\mathfrak {q} _ {1} ^ {[ i ]} + \mathfrak {q} _ {2} ^ {[ d - i ]}; \mathbf {M}).
\]

\begin{enumerate}
\def\labelenumi{\alph{enumi})}
\setcounter{enumi}{5}
\tightlist
\item
  证明有：
\end{enumerate}

\[
e \left(\overline {{{\mathfrak {q}}}} _ {1} ^ {[ \alpha_ {1} ]} + \dots + \overline {{{\mathfrak {q}}}} _ {k} ^ {[ \alpha_ {k} ]}; \mathrm{M}\right) = e \left(\mathfrak {q} _ {1} ^ {[ \alpha_ {1} ]} + \dots + \mathfrak {q} _ {k} ^ {[ \alpha_ {k} ]}; \mathrm{M}\right)
\]

其中 \(\overline{\mathfrak{q}}_{i}\) 表示理想 \(\mathfrak{q}_{i}\) 在 A 中的整闭包（参见习题 19）。

\begin{enumerate}
\def\labelenumi{\arabic{enumi})}
\setcounter{enumi}{20}
\tightlist
\item
  设 A 和 A' 为两个诺特环，q（分别为 q'）为 A（分别为 A'）的理想，包含于 A（分别为 A'）的根基中，且 \(\operatorname{long}(A/q)\)（分别为 \(\operatorname{long}(A'/q'))\) 有限。设 Q 为理想 \(q \otimes A' + A \otimes q'\)，它属于环 \(A \otimes_{\mathbf{Z}} A'\)。证明它包含于 \(A \otimes_{\mathbf{Z}} A'\) 的根基中，并且有等式
\end{enumerate}

\[
e _ {\mathfrak {Q}} (\mathrm{A} \otimes_ {\mathbf {Z}} \mathrm{A} ^ {\prime}) = e _ {\mathfrak {q}} (\mathrm{A}). e _ {\mathfrak {q} ^ {\prime}} (\mathrm{A} ^ {\prime}).
\]

22) * 设 A 为诺特整环，且为局部完备环（分别为域 k 上的有限生成代数），\(a \in A\) 为非可逆非零元素。作下列假设：

α) \(\operatorname{Supp}(A/aA)\) 有唯一极小元 p。

\(\beta)\) 有 \(aA_{p} = pA_{p}\)。

γ) A/p 是整闭环。

本习题旨在证明，对于 A 中每个包含 a 的极大理想 m，环 \(A_{m}\) 是整闭的，且有 p = aA（“Hironaka 引理”）。本习题将使用串联性。在第 X 章中将看到，整的诺特局部完备环满足此性质。域上的有限型整代数也满足此性质（§ 2, n° 4, 定理 3）。

\begin{enumerate}
\def\labelenumi{\alph{enumi})}
\item
  证明 \(A_p\) 是离散赋值环（§ 3, n° 1, 命题 1 及 VI, § 3, n° 6, 命题 9）。
\item
  作 \(\alpha\) 和 \(\beta\) 的假设，并假设 A 是整闭的。证明有 \(\mathfrak{p} = a\mathbf{A}\)（VII, § 1, n° 4, 命题 8）。
\item
  设 A' 为 A 的整闭包。证明 A' 是整的诺特局部完备环（分别证明它是 k 上的有限型代数）。（注意 A 是日式环（IX, § 4, n° 2, 定理 2），因而 A' 是整的、诺特的、半局部且完备的，从而是局部的（III, § 2, n° 13, 命题 19 的推论）。当 A 是 k 上的有限型代数时，使用 V, § 3, n° 2, 定理 2。）
\item
  证明有 A' ⊗A A\_p = A\_p（使用 a) 及 V, § 1, n° 5, 命题 16）。推出存在 A' 中位于 p 之上的唯一素理想 p'，且有 A\_p = A'\_p'、aA'\_p' = p'A'\_p'。
\item
  设 q 为 A'/aA' 的支集中的极小元。证明 q 是 A' 中位于 p 之上的素理想。（注意 q 的高度为 1（§ 3, n° 1, 命题 1），因而根据串联性有 \(\dim(A'/q)=\dim(A)-1\)。推出 \(\dim(A/(q \cap A))=\dim(A)-1\)，从而 \(q \cap A = p\)。）
\item
  从 e) 推出 \(aA'\) 是 \(A'\) 的素理想，且有 \(aA' = p' = pA'\)（可以使用 b)）。
\end{enumerate}

g) 证明有 \(\mathbf{A}_{\mathfrak{m}}^{\prime} = \mathbf{A}_{\mathfrak{m}}\) 对 A 中每个包含 \(a\) 的极大理想 m 成立。（由包含关系 \(\mathbf{A} \subset \mathbf{A}' \subset \mathbf{A}_{\mathfrak{p}}\) 推出包含关系 \(\mathbf{A}/\mathfrak{p} \subset \mathbf{A}'/\mathfrak{p}' \subset \mathbf{A}_{\mathfrak{p}}/\mathfrak{p}\mathbf{A}_{\mathfrak{p}}\)。推出 \(\mathbf{A}/\mathfrak{p}\) 和 \(\mathbf{A}'/\mathfrak{p}'\) 具有相同的分式域，从而根据 \(\gamma)\)，有 \(\mathbf{A}/\mathfrak{p} = \mathbf{A}'/\mathfrak{p}' = \mathbf{A}' \otimes_{\mathbf{A}} (\mathbf{A}/\mathfrak{p})\)，最后一个等式由 \(f)\) 得出。因此有 \((\mathbf{A}'/\mathbf{A}) \otimes_{\mathbf{A}} (\mathbf{A}/\mathfrak{p}) = 0\)。得出结论。）然后完成 Hironaka 引理的证明。 *

\begin{enumerate}
\def\labelenumi{\arabic{enumi})}
\setcounter{enumi}{22}
\tightlist
\item
  \begin{enumerate}
  \def\labelenumii{\alph{enumii})}
  \tightlist
  \item
    设 A 为满足 \(e(A) = 1\) 的诺特局部环。证明 A 有唯一的极小素理想 \(p\)，使得 \(\dim(A) = \dim(A/p)\)。此外证明有 \(\text{long}(A_p) = 1\) 及 \(e(A/p) = 1\)（\S 7, no. 1, 注记 4）。
  \end{enumerate}
\end{enumerate}

\begin{enumerate}
\def\labelenumi{\alph{enumi})}
\setcounter{enumi}{1}
\tightlist
\item
  回忆（p. 82, 习题 10），若对于 A 的每个极小素理想 q 都有 \(\dim(A/q) = \dim(A)\)，则称局部环是等维的；若 A-模 A 没有嵌入相伴素理想，则称其没有嵌入素理想（IV, § 2, no. 3, 注记）。* 将在第 X 章看到，Macaulay 局部环的整商的完备化具有这两个性质。* 设 A 为无嵌入素理想的等维诺特局部环，且 \(e(A) = 1\)。证明 A 是整环。
\end{enumerate}

24) * 设 A 为诺特局部环，\(\widehat{A}\) 为其完备化。本习题旨在证明下列性质等价：

\begin{enumerate}
\def\labelenumi{(\roman{enumi})}
\item
  A 是正则环；
\item
  \(\hat{\mathbf{A}}\) 是等维的且没有嵌入素理想，并且 \(e(\mathbf{A}) = 1\)；
\item
  \(\hat{\mathbf{A}}\) 是整环且 \(e(\mathbf{A}) = 1\)。
\end{enumerate}

证明蕴含 (i) \(\Rightarrow\) (ii) \(\Rightarrow\) (iii)。注意 (iii) 蕴含等式 \(e(\hat{\mathbf{A}}) = 1\)，且为证明 (iii) \(\Rightarrow\) (i)，只需证明 \(\hat{\mathbf{A}}\) 是正则环。因此，为证明 (iii) \(\Rightarrow\) (i)，可以假设 A 完备，以下均如此。我们先处理剩余域 \(\kappa_{\mathrm{A}}\) 含有无穷多个元素的情形。对 A 的维数作归纳。考察 \(\dim(A) = 0\) 的情形，以下设 \(\dim(A) > 0\)。此时存在表面元素 \(x \in \mathrm{A}\)（\S 7, no. 5, 注记 4）。根据串联性（习题 22）及 \S 3, no. 1, 命题 1，对于包含 xA 的每个极小素理想 p，都有 \(\dim(A/p) = \dim(A/xA)\)。由于有 \(e(\mathrm{A}/x\mathrm{A}) = 1\)（\S 7, no. 5, 定理 1），包含 xA 的极小素理想中存在唯一一个 p，且有 \(xA_{p} = pA_{p}\)、\(e(\mathrm{A}/p) = 1\)（习题 23）。根据归纳假设，A/p 是正则环，因而是整闭的（\S 5, no. 2, 定理 1 的推论 1）。由 Hironaka 引理（习题 22）推出 p = xA。因此 A/xA 是正则环，并且由于 A 是整环，A 是正则环（\S 5, no. 3, 命题 2）。

现在假设 \(\kappa_{A}\) 是有限域。证明只需证明爆破 A{]}X{[}（IX, 附录, no. 2）是正则环。注意 A 以及由此得到的 A{]}X{[} 同构于正则环的商（IX, § 2, no. 5, 定理 3），且 \(e(A]X[)=1\)。根据 Macaulay 环理论（第 X 章），完备化 \(\widehat{A]X[}\) 是等维的且没有嵌入素理想。推出 \(\widehat{A]X[}\) 是整环（习题 23）且 \(e(\widehat{A]X[})=1\)，因而 \(A]X[\) 是正则环。得出结论。*

\begin{enumerate}
\def\labelenumi{\arabic{enumi})}
\setcounter{enumi}{24}
\tightlist
\item
  设 \(k\) 为一个域，A 为 \(k\)-局部代数，即在有限型 \(k\)-代数的一个素理想处局部化所得。证明 A 是正则环当且仅当 A 是整环且 \(e(A) = 1\)。（可以参考习题 24 所用的方法，或者注意到由于 A 是正则环的整商，局部环 \(\hat{A}\) 是等维的且没有嵌入素理想。）
\end{enumerate}

